\begin{document}

۹۷

شاخه های یخ زده و در تاکستان‌های یخ زده که در امتداد چشم‌انداز تپه‌ای جریان داشتند. ماشین از بوی سیگار مارلبرو و چربی مک‌دونالد پر شده بود؛ هر بار که متوقف می‌شد، راه‌اندازی آن سخت‌تر می‌شد.

اما به زودی، همه این‌ها کم‌ترین مشکلات ما بودند. مسیری که روی نقشه به نظر می‌رسید سفری چهل و پنج دقیقه‌ای به اودسا باشد، تقریباً ده ساعت به درازا کشید.

لحظاتی در زندگی هستند که به اندازه داشتن ماشینی با تانک پر از بنزین، نقشه‌ای از یک قاره کامل که جلوی تو پهن شده، و بهترین هنرمند گریز در دنیا که روی صندلی پشتی اتوبوس تو نشسته است، پر از امکان هستند. حس می‌کنی که می‌توانی به هر جایی که می‌خواهی بروی. به هر حال، مرزها چیستند جز نقاط کنترلی که به تو می‌گویند وارد مرحله جدیدی از ماجراجویی شده‌ای؟ خوب، این ممکن است در بیشتر مواقع درست باشد، اما فرض کن داری در رند مک‌نالی کار می‌کنی و آخرین نسخه نقشه‌ات از اروپای شرقی را به پایان می‌رسانی. و بیایید فرض کنیم که یک کشور کوچک است که مرز با مولداوی دارد — شاید یک دولت کمونیستی شورشی — ولی هیچ دولت دیگری این کشور را، چه دیپلماتیک و چه به طور دیگر، به رسمیت نمی‌شناسد. حالا تو چه کار می‌کنی؟ آیا آن کشور را روی نقشه‌ات درج می‌کنی یا نه؟

من، یک تردست و یک نجیب‌زادهٔ قلابی، در حال رانندگی در سراسر اروپای شرقی بودیم که به طور کاملاً تصادفی جواب این سؤال را پیدا کردیم. تاکنون این رانندگی بی‌حاصل بوده است. میستری زیر پتو در صندلی عقب کج شده بود، نتوانسته بود از بیماری‌اش بیرون بیاید. از چشم انداز دراماتیک رومانیایی برفی که روز به روز عبور می‌کردیم بی‌خبر بود، او چشمانش را با کلاهش پوشانده و شکایت می‌کرد. گاهی وقت‌ها، به هوشیاری برمی‌گشت و محتوای ذهنش را بازگو می‌کرد. و هر بار محتوای ذهنش نقشه‌ای دیگر بود.

"نقشه‌ای که دارم این است که آمریکای شمالی را گردش کنم و برنامه‌هایم را در کلوب‌های شبانه تبلیغ کنم." او گفت. "فقط نیاز دارم که یک توهم خوب برای رقاص‌ها طراحی کنم. تو می‌توانی مشاور من باشی، استایل. تصور کن: تو و من در کلوب‌های شبانه گردش می‌کنیم و همه دخترها را برای تماشای نمایش روز بعد می‌بریم."

چند روز آرام در کیشیناو گذشت که تنها زنان زیبا را روی جلد مجله‌ها و بیلبوردها می‌دیدیم، و فکر کردیم، "چرا همین‌جا متوقف شویم؟" اودسا بسیار نزدیک بود. شاید ماجراجویی که دنبالش می‌گشتیم در آینده‌ای دورتر قرار داشت.

بنابراین جمعه‌ای سرد و برفی کیشیناو را ترک کردیم و به سمت مرز اوکراین به شرق شمالی حرکت کردیم. جاده‌های پوشیده از برف که از شهر بیرون می‌رفت، تنها به واسطهٔ خطوط لاستیکی یخ‌زده در افق شناخته می‌شد. این منظره شبیه یک صحنه از یک رمان عاشقانه حماسی روسی به نظر می‌رسید، با شاخه‌های درختان که با کریستال‌های یخ پوشانده شده‌اند و... 

\end{document}